\documentclass{article}

% Language setting
% Replace `english' with e.g. `spanish' to change the document language
%\usepackage[english,portuguese]{babel}
\usepackage[brazil]{babel}
\usepackage[utf8]{inputenc}

% Set page size and margins
% Replace `letterpaper' with `a4paper' for UK/EU standard size
\usepackage[a4paper,top=2cm,bottom=2cm,left=2cm,right=2cm,marginparwidth=1.75cm]{geometry}

% Useful packages
\usepackage{amsmath,amsthm,amsfonts,float}
\usepackage{graphicx}
\usepackage[colorlinks=true, allcolors=blue]{hyperref}
\usepackage{arydshln}

%\newenvironment{namebox}
%{\begin{flushleft}Nome: \underline{\hspace{6cm}}\end{flushleft}} % Adjust the length of the underline as needed

\newcommand{\na}{\mathbb{N}} % Set of integers
\newcommand{\naz}{\mathbb{N}_0} % Set of integers
\newcommand{\Z}{\mathbb{Z}} % Set of integers
\newcommand{\re}{\mathbb{R}} % Set of real numbers
\newcommand{\sinal}{\mathrm{sinal}} % Signal operador

\newcommand{\vp}{{\mathbf p}}
\newcommand{\vP}{{\mathbf P}}

\newcommand{\vpi}{{\boldsymbol \pi}}
%\newcommand{\Pr}{{\mathbb P}}

\date{--/--/----} % Removendo a data

\begin{document}

\begin{flushleft}
	Universidade Federal da Bahia\\
	Instituto de Matemática e Estatística\\
	Departamento de Estatística\\
	MATF14 - Estatística Econômica I - 2026.1\\
	Professor: Rodney Fonseca \\
%	Data: 24/09/2024
\end{flushleft}

\begin{center}
	\begin{Large}
		Segunda chamada da prova 2 {\color{blue} - GABARITO} \\
	\end{Large}	
%	Data: 14/05/2025
\end{center}

\fbox{\begin{minipage}{30em}
		Nome:\\
		\\
\end{minipage}}
\hspace{2em} Nota: 

\vspace{1em}

\begin{flushleft}
\begin{small}

\begin{itemize}
	\item Todas as respostas devem estar escritas à caneta azul ou preta. Mantenha a sua letra legível.
	\item Escreva o seu nome em todas as folhas com respostas.
	\item Justifique as suas respostas e apresente os devidos cálculos.
\end{itemize}

\end{small}
\end{flushleft}

\vspace{1em}

% va discreta simples
% identificar binomial
% poisson parecida com lista
% va continua identificar constante e calcular prob
% questao de normal

\begin{enumerate}

% questao de achar a constante feita em sala
\item Seja $\lambda > 0$ uma constante e seja $X$ uma variável aleatória cuja função densidade de probabilidade é 
\begin{align*}
f(x) = \begin{cases}
C \cdot e^{-\lambda x}, & x \geq 0, \\
0, & x < 0.
\end{cases}
\end{align*}
\begin{enumerate}
	\item (1 ponto) Encontre o valor de $C$ para $f(x)$ ser uma legítima função densidade de probabilidade.\\
	{\color{blue}
	Como $f(x)$ deve ser não-negativa, devemos ter $C\geq 0$. Além disso, $\int_{-\infty}^{\infty} f(x)dx = 1$. Note que
	\begin{align*}
	\int_{-\infty}^{\infty} f(x)dx = \int_{0}^{\infty} C \cdot e^{-\lambda x}dx
	= C \left( -\frac{e^{-\lambda x}}{\lambda} \Big|_{0}^{\infty} \right) = \frac{C}{\lambda}.
	\end{align*}
	Logo, temos que $C = \lambda$.
	}

	\item (1 ponto) Calcule a função de distribuição acumulada de $X$.\\
	{\color{blue}
	Como $f(x)=0$ para $x<0$, então a função de distribuição acumulada $F$ vale $F(t)=0$ para $t < 0$. Seja $t \geq 0$, então
	\begin{align*}
	F(t) = P(X\leq t) = \int_{0}^{\infty} \lambda \cdot e^{-\lambda x}dx
	=  -e^{-\lambda x} \Big|_{0}^{t} = 1 - e^{\lambda t} .
	\end{align*}
	Logo,
	\begin{align*}
	F(t) = \begin{cases}
	0, & t<0 \ \\
	1 - e^{\lambda t}, & t\geq 0 .
	\end{cases}
	\end{align*}
	}
	
	\item (1 ponto) Assuma que $\lambda = 1/5$. Calcule $P(X \leq 10)$.\\
	{\color{blue}
	Segue da função de distribuição acumulada que
	$$
	P(X \leq 10) = F(10) = 1 - e^{-\frac{1}{5} 10} = 1 - e^{-2} \approx 0,86.
	$$
	}
\end{enumerate}

% 
\item Considere uma variável aleatória discreta $X$ cuja distribuição de probabilidade é apresentada a seguir:
    \begin{center}
    \begin{tabular}{c|cccccc}
    \hline
    $k$ & 1 & 2 & 3 & 4 & 5 \\
    \hline
    $P(X=k)$ & 1/10 & 1/10 & 4/10 & 3/10 & 1/10 \\
    \hline
    \end{tabular}\\
    \end{center}
    \begin{enumerate}
    	\item (1 ponto) Calcule a probabilidade $P(|X-3|>1)$.\\
	{\color{blue}
	Note que $|X-3|<1$ se $X>4$ ou $X<2$. Logo,
	$$
	P(|X-3|>1) = P(X>4) + P(X<2) = \frac{1}{10} + \frac{1}{10} = \frac{1}{5}.
	$$
	}
	
    	\item (1 ponto) Calcule o desvio padrão de $X$.\\
	{\color{blue}
	A esperança de $X$ é
	$$
	E(X) = 1\cdot\frac{1}{10} + 2\cdot\frac{1}{10} + 3\cdot\frac{4}{10} + 4\cdot\frac{3}{10} + 5\cdot\frac{1}{10}
	= \frac{1 + 2 +12 + 12 + 5}{10} = \frac{32}{10} = 3,2.
	$$
	A esperança de $X^2$ é
	$$
	E(X^2) = 1\cdot\frac{1}{10} + 4\cdot\frac{1}{10} + 9\cdot\frac{4}{10} + 16\cdot\frac{3}{10} + 25\cdot\frac{1}{10}
	= \frac{1 + 4 + 36 + 48 + 25}{10} = \frac{114}{10} = 11,4.
	$$
	A variância de $X$ é $Var(X) = 11,4 - 3,2^2 = 11,4 - 10,24 = 1,16$. Logo, o desvio padrão de $X$ é $dp(X) = 1,077$.
	}
	
    	\item (1 ponto) Calcule a função de distribuição acumulada de $X$.\\
	{\color{blue}
	A função de distribuição acumulada de $X$ é:
	$$
	F(t) = \begin{cases}
	0, & t<1, \\
	1/10, & 1\leq t<2, \\
	2/10, & 2\leq t<3, \\
	6/10, & 3\leq t<4, \\
	9/10, & 4\leq t<5, \\
	1, & t \geq 5 .
	\end{cases}
	$$
	}
    \end{enumerate}

% questao de va discreta feita em sala (prob, media e variancia)
\item Uma central de atendimento recebe em média 2 chamadas por minuto. Seja $X$ a variável aleatória representando o número de chamadas recebidas por minuto e assuma que $X$ tem distribuição Poisson. 
\begin{enumerate}
	\item (1 ponto) Calcule a probabilidade da central não receber chamada alguma durante um minuto.\\
	{\color{blue}
	Note que $X \sim Pois(2)$. Logo, a probabilidade da central não receber chamada alguma durante um minuto é $P(X=0) = e^{-2} \approx 0,13$.
	}
	
	\item (1 ponto) Calcule a probabilidade do número de chamadas ficar acima da média.\\
	{\color{blue}
	Temos que a média é $E(X) = 2$. Logo, a probabilidade do número de chamadas ficar acima da média é
	\begin{align*}
	P(X > E(X)) &= P(X>2) = 1 - P(X\leq 2)
	= 1 - P(X=0) - P(X=1) - P(X=2)\\
	&= 1 - e^{-2} - 2e^{-2} - 2e^{-2} \approx 0,32.
	\end{align*}
	}
\end{enumerate}

% questao de normal
% Levin Q 27, pag 152
\item Considere que o retorno anual de um certo fundo de investimentos é normalmente distribuído com média de 8\% e desvio padrão de 20\%.
\begin{enumerate}
  	\item (1 ponto) Calcule a probabilidade do retorno em um ano qualquer ser maior que 10\%.\\
	{\color{blue}
	Seja $X$ a variável aleatória representando o retorno anual do fundo de investimento. Temos que $X\sim N(8, 20^2)$. Logo, a probabilidade do retorno ficar acima de 10\% é:
	$$
	P(X > 10) = 1 - P(X\leq 10) = 1 - P\left( Z \leq \frac{10-8}{20} \right)
	= 1 - P(Z \leq 0,1)
	= 1 - 0,5398
	= 0,4602.
	$$
	}
	
   	\item (1 ponto) Calcule a probabilidade do retorno em um ano qualquer ser menor que $-5$\%.\\
	{\color{blue}
	A probabilidade do retorno ficar abaixo de -5\% é
	$$
	P(X<-5) = P\left( Z \leq \frac{-5-8}{20} \right)
	= P(Z \leq -0,65)
	= 0,2578.
	$$
	}
\end{enumerate} 


    
\end{enumerate}



\vspace{2cm}

\newpage
\begin{center}
	\Large{\textbf{Formulário}}
\end{center}

\begin{small}
\begin{itemize}
	\item A média aritmética é $\bar{x} = \frac{1}{n} \sum_{i=1}^{n} x_i$. Para dados agregados, $\bar{x} = \frac{1}{n} \sum_{i=1}^{k} n_i * x_i$, em que $n_i$ é a frequência absoluta de $x_i$ e $n=\sum_{i=1}^k n_i$.
	\item Com valores ordenados $x_{(1)} \leq \cdots \leq x_{(n)}$, a mediana é $x_{((n+1)/2)}$ se $n$ for ímpar e $\frac{x_{(n/2)} + x_{((n/2) + 1)}}{2}$ se $n$ for par.
	\item A variância é $Var(X) = \frac{1}{n} \sum_{i=1}^n (x_i - \bar{x})^2$. Para dados agregados, $Var(X) = \frac{1}{n} \sum_{i=1}^k n_i * (x_i - \bar{x})^2$, em que $n_i$ é a freq. absoluta de $x_i$ e $n=\sum_{i=1}^k n_i$. Desvio padrão é $dp(X) = \sqrt{Var(X)}$ e o coeficiente de variação é $dp(X)/\bar{x}$.
%	\item O coeficiente de Yule é $Q = (n_{11}n_{22} - n_{12}n_{21})/(n_{11}n_{22} + n_{12}n_{21})$, em que $n_{ij}$ são frequências absolutas da tabela de dupla entrada correspondente.
	\item O coeficiente de correlação é $cor(X,Y) = \frac{1}{n} \sum_{i=1}^n \left( \frac{x_i - \bar{x}}{dp(X)} \right) \left( \frac{y_i - \bar{y}}{dp(Y)} \right)$.
	
	\item O espaço amostral $S$ é o conjunto de todos resultados possíveis de um experimento aleatório. Um evento $A$ é qualquer sub-conjunto desse espaço amostral $S$.
	\item Probabilidade: $P(A) = \frac{\text{número de elementos de $A$}}{\text{número de elementos do espaço amostral}}$
	\item Probabilidade complementar: $P(A^c) = 1 - P(A)$
	\item Se $A$ e $B$ são eventos disjuntos, ou seja, $A\cap B = \emptyset$, então $P(A \cup B) = P(A) + P(B)$.
	\item Probabilidade condicional: se $P(A)>0$, então $P(B|A) = \frac{P(B \cap A)}{P(A)}$.
	\item Independência: se $A$ e $B$ são eventos independentes, então $P(A \cap B) = P(A) P(B)$.
%	\item Se $A$ e $B$ são eventos independentes e $P(A)>0$, então $P(B | A) = P(B)$.
	\item Teorema da probabilidade total: se $\{B_i\}_{i=1}^n$ é partição de $S$, então $P(A) = \sum_{i=1}^n P(A|B_i) P(B_i)$
	\item Teorema de Bayes: se $\{B_i\}_{i=1}^n$ é partição de $S$ e $P(A)>0$, então $P(B_i | A) = \frac{P(A|B_i) P(B_i)}{\sum_{j=1}^n P(A|B_j) P(B_j)} $
\item Se $X$ é uma variável aleatória contínua, então $E(X) = \int x \cdot f(x)dx$. Se $X$ for uma variável aleatória discreta, $E(X) = \sum_k  k \cdot P(X=k)$.  Além disso, $Var(X) = E(X^2) - E(X)^2$ e $dp(X) = \sqrt{Var(X)}$.
\item Se $X\sim Bernoulli(p)$, então $P(X=k) = p^k (1 - p)^{1 - k}$ com $k\in\{0, 1\}$. $E(X) = p$ e $Var(X) = p(1 - p)$.
\item Se $X\sim Bin(n, p)$, então $P(X=k) = \binom{n}{k} p^k (1 - p)^{n - k}$ com $k\in\{0, 1, \ldots, n\}$. $E(X) = n p$ e $Var(X) = n p(1 - p)$.
\item Se $X\sim Pois(\lambda)$, então $P(X=k)=e^{-\lambda}\lambda^k/k!$, com $k\in\{0,1,2,\ldots\}$. $E(X)=Var(X)=\lambda$.
\item Se $X\sim Exp(\lambda)$, então $f(x) = \lambda e^{-\lambda x}$ para $x>0$, $E(X) = 1/\lambda$ e $Var(X) = 1/\lambda^2$.
\item Se $X\sim N(\mu, \sigma^2)$, então $E(X) = \mu$ e $Var(X) = \sigma^2$. Em particular, $Z = (X - \mu)/\sigma \sim N(0,1)$.
\item Se $X\sim Unif(\alpha, \beta)$ com $\beta>\alpha$, então $f(x)=1/(\beta - \alpha)$ para $\alpha\leq x \leq \beta$. $E(X) = (\beta + \alpha)/2$ e $Var(X) = (\beta - \alpha)^2/12$.
\item $\int x^n dx = \frac{x^{n+1}}{n+1}$, $n\neq -1$ e $\int \frac{1}{x} dx = \ln x$.
\item $\int e^{ax} dx = \frac{1}{a} e^{ax}$.	
\end{itemize}
\end{small}

\newpage
\begin{center}
	\noindent\rule{15cm}{0.4pt}
	\Large{\textbf{Rascunho}}
\end{center}




\end{document}




\end{enumerate}


